\begin{proposition}[末态后选择的成功概率与状态误差]
令 $\sigma=U\rho U^\dagger$ 为归一化联合末态，$P=P_{\rm keep}$ 为正交投影，
$q=\operatorname{Tr}[(I-P)\sigma]<1$。保留事件的成功概率为 $s=1-q$，
其条件态为 $\sigma_{\rm c}=P\sigma P/s$。定义迹距离
$D(\sigma,\tau)=\tfrac12\|\sigma-\tau\|_1$，则
\[
D(\sigma,\sigma_{\rm c})\le\sqrt q,\qquad
|\operatorname{Tr}(M\sigma)-\operatorname{Tr}(M\sigma_{\rm c})|
\le\sqrt q\quad(0\le M\le I).
\]
若入射态支撑在 $\Pi_{\rm in}$ 内，便可用
$q\le\epsilon_{\rm proj}(T)$ 给出统一上界。
对电子或环境再取偏迹后，迹距离界仍成立。
\end{proposition}
\begin{proof}
取 $\sigma$ 的一个归一化纯化 $|\Psi\rangle$，辅助纯化空间记为 $R$。
投影成功后的归一化纯化为
\[
|\Phi\rangle=\frac{(P\otimes I_R)|\Psi\rangle}{\sqrt s},\qquad
\langle\Psi|\Phi\rangle=\sqrt s.
\]
两个纯态之差只作用于它们张成的二维空间。
将第二个态写成 $\sqrt s|\Psi\rangle+\sqrt q|\Psi_\perp\rangle$，
可见差矩阵的迹为零、行列式为 $-q$，故其两个非零本征值为 $\pm\sqrt q$。
两纯态的迹距离于是等于 $\sqrt q$。
偏迹不会增加迹距离：对任意效应 $M$，约化后的概率差等于联合空间中
$M\otimes I_R$ 的概率差；对效应取上确界便得这一收缩性。
先对 $R$ 求偏迹，再按需对未观测系统求偏迹，即得所给状态界。
最后，迹为零的厄米差算符的正、负部分具有相同迹，
所以任意 $0\le M\le I$ 的概率差不超过其迹距离。
\end{proof}

这个界比较完整末态与对完整末态作后选择得到的条件态。
它没有把完整 $U$ 换成截断 Hamiltonian 的传播子。
电子可以暂时离开保留空间，再带着相位返回；
出射时看不到泄漏，并不等于中途经过的通道没有作用。

\begin{exbox}[零末态泄漏仍会改变干涉相位]
取三个正交联合通道 $|0\rangle,|1\rangle,|a\rangle$，
保留 $P=|0\rangle\langle0|+|1\rangle\langle1|$，删去 $|a\rangle$。
在旋转绘景中设
\[
H=\hbar g\bigl(|1\rangle\langle a|+|a\rangle\langle1|\bigr),\qquad g>0.
\]
以 $|+\rangle=(|0\rangle+|1\rangle)/\sqrt2$ 入射，
求 $T=\pi/g$ 时的泄漏和电子干涉读出，
并与直接用 $PHP$ 演化的预测比较。

\textbf{解。} $|0\rangle$ 不参与耦合；在 $|1\rangle,|a\rangle$ 子空间中，
$H/\hbar=g\sigma_x$，所以
\[
U(t)|1\rangle=\cos(gt)|1\rangle-\ii\sin(gt)|a\rangle.
\]
任意保留输入的最坏泄漏为 $\epsilon_{\rm proj}(t)=\sin^2(gt)$。
在所给末时刻，它等于零，却有
\[
U(T)|+\rangle=\frac{|0\rangle-|1\rangle}{\sqrt2}=| -\rangle.
\]
完整模型对效应 $M_+=|+\rangle\langle+|$ 给出概率零。
另一方面，$PHP=0$，其传播子在保留空间中为恒等算符，
预测同一效应的概率为一。误差达到最大值，尽管末态泄漏完全消失。

差别来自耦合通道的一次完整往返。
这里两个保留分量之间的相对相位可由干涉读出，不能当作公共相位删去。
若要删去 $|a\rangle$，需要保留它产生的有效能移与相位，
或直接比较完整传播子与候选有效传播子；后一比较由下面的 Duhamel 恒等式控制。
\end{exbox}
